\begin{thebibliography}{45}
\bibitem[1]{ref1} Luis Filipe Coelho Antunes and Lance Fortnow. Worst-case running times for average-case algorithms. In Proceedings of the 24th Annual IEEE Conference on Computational Complexity, CCC 2009, Paris, France, 15-18 July 2009, pages 298–303. IEEE Computer Society, 2009. doi:10.1109/CCC.2009.12.
\bibitem[2]{ref2} Shai Ben-David, Benny Chor, Oded Goldreich, and Michael Luby. On the theory of average case complexity. J. Comput. Syst. Sci., 44(2):193–219, 1992. doi:10.1016/0022-0000(92)90019-F.
\bibitem[3]{ref3} Andrej Bogdanov and Luca Trevisan. Average-case complexity. Found. Trends Theor. Comput. Sci., 2(1), 2006. doi:10.1561/0400000004.
\bibitem[4]{ref4} Andrej Bogdanov and Luca Trevisan. On worst-case to average-case reductions for NP problems. SIAM J. Comput., 36(4):1119–1159, 2006. doi:10.1137/S0097539705446974.
\bibitem[5]{ref5} Harry Buhrman, Lance Fortnow, and Aduri Pavan. Some results on derandomization. Theory Comput. Syst., 38(2):211–227, 2005. doi:10.1007/s00224-004-1194-y.
\bibitem[6]{ref6} Lijie Chen, Shuichi Hirahara, and Neekon Vafa. Average-case hardness of NP and PH from worst-case fine-grained assumptions. In Mark Braverman, editor, 13th Innovations in Theoretical Computer Science Conference, ITCS 2022, January 31 - February 3, 2022, Berkeley, CA, USA, volume 215 of LIPIcs, pages 45:1–45:16. Schloss Dagstuhl - Leibniz-Zentrum für Informatik, 2022. doi:10.4230/LIPIcs.ITCS.2022.45.
\bibitem[7]{ref7} Halley Goldberg and Valentine Kabanets. A simpler proof of the worst-case to average-case reduction for polynomial hierarchy via symmetry of information. Electron. Colloquium Comput. Complex., 7:1–14, 2022. URL: https://eccc.weizmann.ac.il/report/2022/007/.
\bibitem[8]{ref8} Oded Goldreich. The Foundations of Cryptography - Volume 1: Basic Techniques. Cambridge University Press, 2001. doi:10.1017/CBO9780511546891.
\bibitem[9]{ref9} Oded Goldreich and Leonid A. Levin. A hard-core predicate for all one-way functions. In Proceedings of the Twenty-First Annual ACM Symposium on Theory of Computing, STOC ’89, pages 25–32, New York, NY, USA, 1989. Association for Computing Machinery. doi:10.1145/73007.73010.
\bibitem[10]{ref10} Shuichi Hirahara. Characterizing average-case complexity of PH by worst-case meta-complexity. In Sandy Irani, editor, 61st IEEE Annual Symposium on Foundations of Computer Science, FOCS 2020, Durham, NC, USA, November 16-19, 2020, pages 50–60. IEEE, 2020. doi:10.1109/FOCS46700.2020.00014.
\bibitem[11]{ref11} Shuichi Hirahara. Average-case hardness of NP from exponential worst-case hardness assump- tions. In Samir Khuller and Virginia Vassilevska Williams, editors, STOC ’21: 53rd Annual ACM SIGACT Symposium on Theory of Computing, Virtual Event, Italy, June 21-25, 2021, pages 292–302. ACM, 2021. doi:10.1145/3406325.3451065.
\bibitem[12]{ref12} Shuichi Hirahara. Symmetry of information in heuristica. Manuscript, 2021.
\bibitem[13]{ref13} Shuichi Hirahara and Mikito Nanashima. On worst-case learning in relativized heuristica. In Symposium on Foundations of Computer Science (FOCS), 2021.
\bibitem[14]{ref14} Shuichi Hirahara and Rahul Santhanam. Errorless versus error-prone average-case complexity. In Mark Braverman, editor, 13th Innovations in Theoretical Computer Science Conference, ITCS 2022, January 31 - February 3, 2022, Berkeley, CA, USA, volume 215 of LIPIcs, pages 84:1–84:23. Schloss Dagstuhl - Leibniz-Zentrum für Informatik, 2022. doi:10.4230/LIPIcs.ITCS.2022.84.
\bibitem[15]{ref15} Shuichi Hirahara and Rahul Santhanam. Excluding PH pessiland. In Mark Braverman, editor, 13th Innovations in Theoretical Computer Science Conference, ITCS 2022, January 31 - February 3, 2022, Berkeley, CA, USA, volume 215 of LIPIcs, pages 85:1–85:25. Schloss Dagstuhl - Leibniz-Zentrum für Informatik, 2022. doi:10.4230/LIPIcs.ITCS.2022.85.
\bibitem[16]{ref16} Russell Impagliazzo. A personal view of average-case complexity. In Proceedings of the Tenth Annual Structure in Complexity Theory Conference, Minneapolis, Minnesota, USA, June 19-22, 1995, pages 134–147, 1995. doi:10.1109/SCT.1995.514853.
\bibitem[17]{ref17} Russell Impagliazzo. Relativized separations of worst-case and average-case complexities for NP. In Proceedings of the 26th Annual IEEE Conference on Computational Complexity, CCC 2011, San Jose, California, USA, June 8-10, 2011, pages 104–114, 2011. doi:10.1109/CCC.2011.34.
\bibitem[18]{ref18} Russell Impagliazzo and Leonid A. Levin. No better ways to generate hard NP instances than picking uniformly at random. In 31st Annual Symposium on Foundations of Computer Science, St. Louis, Missouri, USA, October 22-24, 1990, Volume II, pages 812–821. IEEE Computer Society, 1990. doi:10.1109/FSCS.1990.89604.
\bibitem[19]{ref19} Russell Impagliazzo and Avi Wigderson. P = BPP if E requires exponential circuits: De- randomizing the XOR lemma. In Frank Thomson Leighton and Peter W. Shor, editors, Proceedings of the Twenty-Ninth Annual ACM Symposium on the Theory of Computing, El Paso, Texas, USA, May 4-6, 1997, pages 220–229. ACM, 1997. doi:10.1145/258533.258590.
\bibitem[20]{ref20} Michael J. Kearns, Robert E. Schapire, and Linda Sellie. Toward efficient agnostic learning. Mach. Learn., 17(2-3):115–141, 1994. doi:10.1007/BF00993468.
\bibitem[21]{ref21} Pravesh K. Kothari and Roi Livni. Improper learning by refuting. In Anna R. Karlin, editor, 9th Innovations in Theoretical Computer Science Conference, ITCS 2018, January 11-14, 2018, Cambridge, MA, USA, volume 94 of LIPIcs, pages 55:1–55:10. Schloss Dagstuhl - Leibniz-Zentrum für Informatik, 2018. doi:10.4230/LIPIcs.ITCS.2018.55.
\bibitem[22]{ref22} Clemens Lautemann. BPP and the polynomial hierarchy. Inf. Process. Lett., 17(4):215–217, 1983. doi:10.1016/0020-0190(83)90044-3.
\bibitem[23]{ref23} Troy Lee. Kolmogorov complexity and formula lower bounds. PhD thesis, University of Amsterdam, 2006.
\bibitem[24]{ref24} Leonid A. Levin. Laws of information conservation (nongrowth) and aspects of the foundation of probability theory. Problemy Peredachi Informatsii, 10(3):30–35, 1974.
\bibitem[25]{ref25} Leonid A. Levin. Average case complete problems. SIAM J. Comput., 15(1):285–286, 1986. doi:10.1137/0215020.
\bibitem[26]{ref26} Ming Li and Paul M. B. Vitányi. An introduction to Kolmogorov complexity and its applications. Springer, 2019.
\bibitem[27]{ref27} Ming Li and Paul M. B. Vitányi. An Introduction to Kolmogorov Complexity and Its Applications, 4th Edition. Texts in Computer Science. Springer, 2019. doi:10.1007/978-3-030-11298-1.
\bibitem[28]{ref28} Yanyi Liu and Rafael Pass. On one-way functions and kolmogorov complexity. In 61st IEEE Annual Symposium on Foundations of Computer Science, FOCS 2020, Durham, NC, USA, November 16-19, 2020, pages 1243–1254, 2020. doi:10.1109/FOCS46700.2020.00118.
\bibitem[29]{ref29} Yanyi Liu and Rafael Pass. On one-way functions from np-complete problems. Electron. Colloquium Comput. Complex., page 59, 2021. URL: https://eccc.weizmann.ac.il/report/2021/059.
\bibitem[30]{ref30} Yanyi Liu and Rafael Pass. On the possibility of basing cryptography on $\mathrm{EXP}\neq\mathrm{BPP}$. In Advances in Cryptology - CRYPTO 2021 - 41st Annual International Cryptology Conference, CRYPTO 2021, Virtual Event, August 16-20, 2021, Proceedings, Part I, pages 11–40, 2021. doi:10.1007/978-3-030-84242-0\_2.
\bibitem[31]{ref31} Luc Longpré and Osamu Watanabe. On symmetry of information and polynomial time invertibility. Inf. Comput., 121(1):14–22, 1995. doi:10.1006/inco.1995.1120.
\bibitem[32]{ref32} Zhenjian Lu and Igor C. Oliveira. An efficient coding theorem via probabilistic representations and its applications. In Nikhil Bansal, Emanuela Merelli, and James Worrell, editors, 48th International Colloquium on Automata, Languages, and Programming, ICALP 2021, July 12-16, 2021, Glasgow, Scotland (Virtual Conference), volume 198 of LIPIcs, pages 94:1–94:20. Schloss Dagstuhl - Leibniz-Zentrum für Informatik, 2021. doi:10.4230/LIPIcs.ICALP.2021.94.
\bibitem[33]{ref33} Zhenjian Lu, Igor C. Oliveira, and Rahul Santhanam. Pseudodeterministic algorithms and the structure of probabilistic time. In Samir Khuller and Virginia Vassilevska Williams, editors, STOC ’21: 53rd Annual ACM SIGACT Symposium on Theory of Computing, Virtual Event, Italy, June 21-25, 2021, pages 303–316. ACM, 2021. doi:10.1145/3406325.3451085.
\bibitem[34]{ref34} Zhenjian Lu, Igor C. Oliveira, and Marius Zimand. Optimal coding theorems in time-bounded kolmogorov complexity. In 49th International Colloquium on Automata, Languages and Programming, 2022.
\bibitem[35]{ref35} Peter Bro Miltersen, N. V. Vinodchandran, and Osamu Watanabe. Super-polynomial versus half-exponential circuit size in the exponential hierarchy. In Takao Asano, Hiroshi Imai, D. T. Lee, Shin-Ichi Nakano, and Takeshi Tokuyama, editors, Computing and Combinatorics, 5th Annual International Conference, COCOON ’99, Tokyo, Japan, July 26-28, 1999, Proceedings, volume 1627 of Lecture Notes in Computer Science, pages 210–220. Springer, 1999. doi:10.1007/3-540-48686-0\_21.
\bibitem[36]{ref36} Ashish V. Naik, Kenneth W. Regan, and D. Sivakumar. On quasilinear-time complexity theory. Theor. Comput. Sci., 148(2):325–349, 1995. doi:10.1016/0304-3975(95)00031-Q.
\bibitem[37]{ref37} Noam Nisan and Avi Wigderson. Hardness vs randomness. J. Comput. Syst. Sci., 49(2):149–167, 1994. doi:10.1016/S0022-0000(05)80043-1.
\bibitem[38]{ref38} Igor C. Oliveira. Randomness and intractability in kolmogorov complexity. In Christel Baier, Ioannis Chatzigiannakis, Paola Flocchini, and Stefano Leonardi, editors, 46th International Colloquium on Automata, Languages, and Programming, ICALP 2019, July 9-12, 2019, Patras, Greece, volume 132 of LIPIcs, pages 32:1–32:14. Schloss Dagstuhl - Leibniz-Zentrum für Informatik, 2019. doi:10.4230/LIPIcs.ICALP.2019.32.
\bibitem[39]{ref39} Igor C. Oliveira and Rahul Santhanam. Pseudodeterministic constructions in subexponential time. In Hamed Hatami, Pierre McKenzie, and Valerie King, editors, Proceedings of the 49th Annual ACM SIGACT Symposium on Theory of Computing, STOC 2017, Montreal, QC, Canada, June 19-23, 2017, pages 665–677. ACM, 2017. doi:10.1145/3055399.3055500.
\bibitem[40]{ref40} Hanlin Ren and Rahul Santhanam. Hardness of KT characterizes parallel cryptography. In 36th Computational Complexity Conference, CCC 2021, July 20-23, 2021, Toronto, Ontario, Canada (Virtual Conference), pages 35:1–35:58, 2021. doi:10.4230/LIPIcs.CCC.2021.35.
\bibitem[41]{ref41} Shai Shalev-Shwartz and Shai Ben-David. Understanding Machine Learning - From Theory to Algorithms. Cambridge University Press, 2014. URL: http://www.cambridge.org/de/academic/subjects/computer-science/pattern-recognition-and-machine-learning/ understanding-machine-learning-theory-algorithms.
\bibitem[42]{ref42} Ronen Shaltiel and Christopher Umans. Simple extractors for all min-entropies and a new pseudorandom generator. J. ACM, 52(2):172–216, 2005. doi:10.1145/1059513.1059516.
\bibitem[43]{ref43} Luca Trevisan and Salil P. Vadhan. Pseudorandomness and average-case complexity via uniform reductions. Computational Complexity, 16(4):331–364, 2007.
\bibitem[44]{ref44} Salil P. Vadhan. On learning vs. refutation. In Satyen Kale and Ohad Shamir, editors, Proceed- ings of the 30th Conference on Learning Theory, COLT 2017, Amsterdam, The Netherlands, 7-10 July 2017, volume 65 of Proceedings of Machine Learning Research, pages 1835–1848. PMLR, 2017. URL: http://proceedings.mlr.press/v65/vadhan17a.html.
\bibitem[45]{ref45} Emanuele Viola. On constructing parallel pseudorandom generators from one-way functions. In 20th Annual IEEE Conference on Computational Complexity (CCC 2005), 11-15 June 2005, San Jose, CA, USA, pages 183–197, 2005. doi:10.1109/CCC.2005.16.
\end{thebibliography}
